\documentclass[12pt]{article}
\usepackage{amsmath,amssymb,amsfonts,amsthm}
\usepackage[margin=1in]{geometry}

\begin{document}

\begin{center}
{\Large\bfseries Chapter 4, Problem 12}\\[6pt]
Byron \& Fuller, \textit{Mathematics of Classical and Quantum Physics}
\end{center}

\bigskip

\textbf{Problem.} The Lorentz group.
\begin{itemize}
\item[(a)] Evaluate the determinant of the Lorentz matrix. Comments?
\item[(b)] Prove that the set of Lorentz transformations in one velocity direction form a group.
\end{itemize}

\bigskip
\textbf{Solution.}

\textbf{(a)} The Lorentz transformation for a boost in the $x$-direction with velocity $v$ is given by the matrix (using $ct, x, y, z$ coordinates):
\[
L = \begin{pmatrix} \gamma & -\beta\gamma & 0 & 0 \\ -\beta\gamma & \gamma & 0 & 0 \\ 0 & 0 & 1 & 0 \\ 0 & 0 & 0 & 1 \end{pmatrix},
\]
where $\beta = v/c$ and $\gamma = 1/\sqrt{1 - \beta^2}$.

The determinant is:
\[
\det L = 1 \cdot 1 \cdot \det\begin{pmatrix} \gamma & -\beta\gamma \\ -\beta\gamma & \gamma \end{pmatrix} = \gamma^2 - \beta^2\gamma^2 = \gamma^2(1 - \beta^2) = 1.
\]

\textbf{Comment:} The determinant of a proper Lorentz transformation is $+1$. This means Lorentz transformations preserve the orientation of spacetime. More generally, all Lorentz transformations satisfy $\det L = \pm 1$ (from $L^T \eta L = \eta$), and the proper Lorentz transformations (connected to the identity) have $\det L = +1$.

\bigskip
\textbf{(b)} We need to verify the four group axioms for Lorentz boosts along one direction.

Using the rapidity parameter $\alpha = \tanh^{-1}\beta$, the Lorentz boost takes the form:
\[
L(\alpha) = \begin{pmatrix} \cosh\alpha & -\sinh\alpha \\ -\sinh\alpha & \cosh\alpha \end{pmatrix}
\]
(restricting to the $ct$-$x$ subspace; the $y,z$ coordinates are unchanged).

\textbf{Closure:} The product of two boosts is:
\[
L(\alpha_1)L(\alpha_2) = \begin{pmatrix} \cosh\alpha_1 & -\sinh\alpha_1 \\ -\sinh\alpha_1 & \cosh\alpha_1 \end{pmatrix}\begin{pmatrix} \cosh\alpha_2 & -\sinh\alpha_2 \\ -\sinh\alpha_2 & \cosh\alpha_2 \end{pmatrix}.
\]
Using hyperbolic identities:
\[
L(\alpha_1)L(\alpha_2) = \begin{pmatrix} \cosh(\alpha_1+\alpha_2) & -\sinh(\alpha_1+\alpha_2) \\ -\sinh(\alpha_1+\alpha_2) & \cosh(\alpha_1+\alpha_2) \end{pmatrix} = L(\alpha_1 + \alpha_2).
\]
This is another Lorentz boost, so the set is closed under composition.

\textbf{Identity:} $L(0) = I$ is the identity element.

\textbf{Inverse:} $L(\alpha)^{-1} = L(-\alpha)$, since $L(\alpha)L(-\alpha) = L(0) = I$.

\textbf{Associativity:} Matrix multiplication is associative.

Therefore the set of Lorentz boosts in one direction forms a group (isomorphic to $(\mathbb{R}, +)$ via the rapidity). \qed

\end{document}
